\documentclass[aspectratio=169]{beamer}
%\usetheme{Ilmenau}
%\usecolortheme{beaver}

\renewcommand{\thefootnote}{\alph{footnote}}
\usepackage[style=numeric, sorting=none, backend=bibtex]{biblatex}
\addbibresource{presentation.bib}
\AtBeginBibliography{\tiny}

% Change the style of the printing biblio to let appear the numbers like in the footnote
\setbeamertemplate{bibliography item}{\insertbiblabel}

\makeatletter
\def\blfootnote{\xdef\@thefnmark{}\@footnotetext}
\makeatother

% Define the new cite command using the footnote
\DeclareCiteCommand{\fcite}
    {\usebibmacro{prenote}} 
    {%
        \renewcommand{\thefootnote}{\arabic{footnote}}% Switch to footnote with numbers.
        \footnotemark[\thefield{labelnumber}]% Add the mark corresponding to the number entry.
        \footnotetext[\thefield{labelnumber}]{%  Add the footnote text with same number entry.
		\usedriver{}{\thefield{entrytype}}}%
%        \blfootnote{%  Add the footnote text with same number entry.
%		$[$\thefield{labelnumber}$]$ \usedriver{}{\thefield{entrytype}}%
%        }%        
%        \renewcommand{\thefootnote}{\alph{footnote}}% Switch back to footnote with letters.
    }
    {{\usebeamercolor[fg]{footnote mark}\textsuperscript,}}
    %{\multicitedelim}
    {\usebibmacro{postnote}} 
    
    
% Define the new cite command using the footnote, use this for previously referenced worsk, so you keep the previous number
\DeclareCiteCommand{\ffcite}
    {\usebibmacro{prenote}} 
    {%
        \renewcommand{\thefootnote}{\arabic{footnote}}% Switch to footnote with numbers.
        \footnotemark[\thefield{labelnumber}]% Add the mark corresponding to the number entry.
%        \renewcommand{\thefootnote}{\alph{footnote}}% Switch back to footnote with letters.
    }
    {{\usebeamercolor[fg]{footnote mark}\textsuperscript,}}
    %{\multicitedelim}
    {\usebibmacro{postnote}}   
   

% Define the new cite command using the footnote for referencing images
\DeclareCiteCommand{\imfcite}
    {\usebibmacro{prenote}} 
    {%
        \renewcommand{\thefootnote}{\arabic{footnote}}% Switch to footnote with numbers.
        %\footnotemark[\thefield{labelnumber}]% Add the mark corresponding to the number entry.
        \blfootnote{%  Add the footnote text with same number entry.
		Image adopted from $[$\thefield{labelnumber}$]$: \usedriver{}{\thefield{entrytype}}%
        }%
        \renewcommand{\thefootnote}{\alph{footnote}}% Switch back to footnote with letters.
    }
    {{\usebeamercolor[fg]{footnote mark}\textsuperscript{,}}%
    }
    %{\multicitedelim}
    {\usebibmacro{postnote}}      
   

\usepackage{bibentry}

\usepackage[american]{babel}
%\usepackage[utf8]{inputenc}
\usepackage{graphicx}
\usepackage{adjustbox}
\usepackage{xcolor}
\usepackage{tikz}
\usepackage{wrapfig}
\usepackage{multirow}
\usepackage{appendixnumberbeamer}
\usepackage{contour}
\usepackage{listings}
\usepackage{pgfplots}
\usepackage{array}
\usepackage{tikz}
\usepgfplotslibrary{colormaps}

%\usepackage{DejaVuSans}
\usepackage[scaled]{berasans}
\renewcommand*\familydefault{\sfdefault}  %% Only if the base font of the document is to be sans serif
\usepackage[T1]{fontenc}

    

% Make tables nice otherwise the text will not align vertically with the table
 \newsavebox\MBox
\newcommand\Cline[2][red]{{\sbox\MBox{$#2$}%
  \rlap{\usebox\MBox}\color{#1}\rule[-2.2\dp\MBox]{\wd\MBox}{1pt}}}
  
\newcommand{\secframe}[2]{
  \setbeamercolor{background canvas}{bg=DefaultColor}
\begin{frame}[plain, noframenumbering]
\frametitle{~}
  \color{white}
  \Large #1 \\
  \huge #2
\end{frame}
\setbeamercolor{background canvas}{bg=}
}


\usepackage{bm}

\renewcommand{\vec}[1]{\bm{#1}}
\newcommand{\mat}[1]{\mathbf{#1}}

\newcommand{\norm}[1]{\left\lVert#1\right\rVert}
\DeclareMathOperator*{\argmax}{arg\,max}


% Uncomment these for FMFI orange colors
%\definecolor{DefaultColor}{HTML}{ff8200}
%\definecolor{Dark10}{HTML}{e66900}
%\definecolor{Dark20}{HTML}{cc4f00}
%\definecolor{Dark30}{HTML}{b33600}
%\definecolor{Dark40}{HTML}{991c00}
%\definecolor{Dark50}{HTML}{800300}
%\definecolor{Dark60}{HTML}{660000}
%\definecolor{Dark70}{HTML}{4d0000}
%\definecolor{Dark80}{HTML}{330000}
%\definecolor{Dark90}{HTML}{1a0000}
%
%\definecolor{Light10}{HTML}{ff9c1a}
%\definecolor{Light20}{HTML}{ffb533}
%\definecolor{Light30}{HTML}{ffcf4d}
%\definecolor{Light40}{HTML}{ffe866}
%\definecolor{Light50}{HTML}{ffff7f}
%\definecolor{Light60}{HTML}{ffff99}
%\definecolor{Light70}{HTML}{ffffb3}
%\definecolor{Light80}{HTML}{ffffcc}
%\definecolor{Light90}{HTML}{ffffe5}


% These are the red UK colors
\definecolor{DefaultColor}{HTML}{9e2b2f}
\definecolor{Dark10}{HTML}{851216}
\definecolor{Dark20}{HTML}{6b0000}
\definecolor{Dark30}{HTML}{520000}
\definecolor{Dark40}{HTML}{380000}
\definecolor{Dark50}{HTML}{1f0000}
\definecolor{Dark60}{HTML}{050000}


\definecolor{Light10}{HTML}{884549}
\definecolor{Light20}{HTML}{d15e62}
\definecolor{Light30}{HTML}{eb787c}
\definecolor{Light40}{HTML}{ff9195}
\definecolor{Light50}{HTML}{ffaaae}
\definecolor{Light60}{HTML}{ffc4c8}
\definecolor{Light70}{HTML}{ffdee2}
\definecolor{Light80}{HTML}{fff7fb}


% add page number
\setbeamertemplate{footline}[frame number]{}
% remove navigation symbols
\setbeamertemplate{navigation symbols}{}



% if you need to change a color of some element you can follow this
\setbeamertemplate{itemize item}[square]
\setbeamercolor{itemize item}{fg = DefaultColor}
\setbeamercolor{itemize subitem}{fg =Light10}
\setbeamercolor{enumerate item}{fg =DefaultColor}
\setbeamercolor{enumerate subitem}{fg =Light10}

\usetheme{default}
\beamertemplatenavigationsymbolsempty
\setbeamercolor{title}{fg=white, bg=DefaultColor}
\setbeamercolor{title page}{fg=white, bg=DefaultColor}
\setbeamercolor{frametitle}{fg=white, bg=DefaultColor}
\setbeamercolor{block title}{bg=DefaultColor, fg= white}
\setbeamercolor{block body}{bg =Light70, fg = black}


\setbeamercolor{bibliography item}{fg=Dark10}
\setbeamercolor*{bibliography entry title}{fg=Dark10}
\setbeamercolor*{bibliography entry author}{fg=Dark10}
\setbeamercolor*{bibliography entry location}{fg=Dark10}
\setbeamercolor*{bibliography entry note}{fg=Dark10}

\setbeamercolor{footnote}{fg=Dark10}
\setbeamercolor{footnote mark}{fg=Dark10}
\setbeamerfont{footnote}{size=\tiny}
\renewcommand\footnoterule{}

\setbeamercolor{page number in head/foot}{fg=gray}
\setbeamerfont{footline}{size=\footnotesize}

\setbeamertemplate{frametitle}
{%
    \nointerlineskip
    \begin{beamercolorbox}[wd=\paperwidth]{frametitle}
        \usebeamerfont{frametitle}%
        \adjustbox{raise=3.3mm, margin=5mm 1mm 1mm 1mm}{\insertframetitle}%
        \hfill
        \adjustbox{margin=2mm 2mm 3mm 2mm}{
        % you can add your project logo here see example comment below
        %\includegraphics[height=10mm]{images/logos/TERAIS}\hspace{2mm}
        \includegraphics[height=10mm]{images/logos/UK_logo_white}%
        \hspace{2mm}%
        \includegraphics[height=10mm]{images/logos/FMFI_logo_white}%
        }
    \end{beamercolorbox}
}

% if you want to use sections, but use no subsection in the top panel you can uncomment this
%\useoutertheme[subsection=false]{miniframes}
%\AtBeginSection[]{\subsection{}}

% removes the subsection indicators, comment this if you want to use them
\renewcommand*{\slideentry}[6]{}

%\setbeamercolor*{mini frame}{fg=white,bg=UKRosy}
\setbeamercolor{section in head/foot}{fg=white, bg=Dark10}
\setbeamertemplate{section in head/foot shaded}{\color{Light10}\usebeamertemplate{section in head/foot}}

%\setbeamertemplate{itemize/enumerate body begin}{\normalsize}
%\setbeamertemplate{itemize/enumerate subbody begin}{\normalsize}

% version with the project logos
%\defbeamertemplate*{title page}{customized}[1][]
%{
%        %\adjustbox{margin=2mm 2mm 3mm 2mm}{
%        \includegraphics[height=18mm]{images/logos/UK_logo_white}
%        \hspace{1mm}
%        \includegraphics[height=18mm]{images/logos/FMFI_logo_text_white}
%        \hspace{1mm}
%         \\
%        
%        %}
%	\vspace{5mm}
%	
%        \huge
%        Sample Presentation Title\\    
%        \normalsize     
%        \vspace{3mm}
%        
%        Ing. Viktor Kocur, PhD.\\        
%             
%        \vspace{5mm}
%        
%        \normalsize
%        
%      
%        \includegraphics[height=18mm]{images/logos/TERAIS}
%	  \hspace{1mm}
%	  \includegraphics[height=18mm]{images/logos/EU}
%}

% version without the project logos
\defbeamertemplate*{title page}{customized}[1][]
{
        %\adjustbox{margin=2mm 2mm 3mm 2mm}{
        \includegraphics[height=18mm]{images/logos/UK_logo_white}
        \hspace{1mm}
        \includegraphics[height=18mm]{images/logos/FMFI_logo_text_white}
        \hspace{1mm}
         \\
        
        %}
	\vspace{5mm}
	
        \huge
        Sample Presentation Title\\    
        \normalsize     
        \vspace{3mm}
        
        Ing. Viktor Kocur, PhD. \\        
     
% You can add some other information
%        \vspace{1mm}   
%        15.1.2024, Bratislava                        
}

\begin{document}
% You can select the language
\selectlanguage{american}
% \selectlanguage{slovak}

\setbeamercolor{background canvas}{bg=DefaultColor}
\begin{frame}[plain]
  \titlepage
\end{frame}
\setbeamercolor{background canvas}{bg=}

\begin{frame}
\frametitle{Contents}
	It is easier to create the table of contents manually.
	\vspace{1em}
      \begin{itemize}\setlength{\itemsep}{1em}
      
      \item Basic Formatting
      \begin{itemize}
      \item Images
      \item Equations
      \item Lists
      \item Blocks
      \item Columns
      \item Tables
      \end{itemize}
      \item Slide Controls
      \item References
      \end{itemize}
\end{frame}

\secframe{~}{Basic Formatting}

\secframe{You can start}{Sections using the secframe command}

\begin{frame}{Importing Images}
  To import an image, use the \texttt{graphicx} package.

  \centering
  \includegraphics[width=0.5\textwidth]{example-image-a}

  \medskip
  \textbf{Caption:} Caption for example-image-a.
\end{frame}

\begin{frame}{Multiple Images}
  You can include multiple images in a grid or side by side.

  \centering
  \begin{minipage}{0.45\textwidth}
    \includegraphics[width=\textwidth]{example-image-a}

    \medskip
    \textbf{Caption:} Caption for example-image-a.
  \end{minipage}
  \hfill
  \begin{minipage}{0.45\textwidth}
    \includegraphics[width=\textwidth]{example-image-b}

    \medskip
    \textbf{Caption:} Caption for example-image-b.
  \end{minipage}
\end{frame}

\begin{frame}{Mathematics in \LaTeX}
  \textbf{Equation Environment:}

  Use the \texttt{equation} environment for numbered equations.

  \begin{equation}
    E = mc^2
  \end{equation}

  This equation is automatically numbered.

  \textbf{Equation* Environment:}

  If you want an unnumbered equation, use the \texttt{equation*} environment.

  \begin{equation*}
    F = ma
  \end{equation*}

  This equation won't have a number.

  \textbf{Inline Mathematics:}

  For inline mathematics, enclose the math expression with dollar signs.

  The famous equation by Euler: $e^{i\pi} + 1 = 0$

  Inline mathematics doesn't get numbered.
\end{frame}

\begin{frame}{Itemization with Two Levels}
  \textbf{Top Level Items:}

  \begin{itemize}
    \item First top-level item.
      \begin{itemize}
        \item First nested item.
        \item Second nested item.
      \end{itemize}
    \item Second top-level item.
      \begin{itemize}
        \item First nested item.
        \item Second nested item.
      \end{itemize}
    \item Third top-level item.
      \begin{itemize}
        \item First nested item.
        \item Second nested item.
      \end{itemize}
  \end{itemize}

  \textbf{Conclusion:}

  Use the \texttt{itemize} environment for easy bullet-point lists, and you can nest items to create a hierarchy.
\end{frame}

\begin{frame}{Enumeration with Two Levels}
  \textbf{Top Level Items:}

  \begin{enumerate}
    \item First top-level item.
      \begin{enumerate}
        \item First nested item.
        \item Second nested item.
      \end{enumerate}
    \item Second top-level item.
      \begin{enumerate}
        \item First nested item.
        \item Second nested item.
      \end{enumerate}
    \item Third top-level item.
      \begin{enumerate}
        \item First nested item.
        \item Second nested item.
      \end{enumerate}
  \end{enumerate}

  \textbf{Conclusion:}

  Use the \texttt{enumerate} environment for numbered lists, and you can nest items to create a hierarchy.
\end{frame}

\begin{frame}{Block Environment Example}
  \begin{block}{Introduction}
    This is a block with a title. It can be used to highlight important information or introduce a section.
  \end{block}

  \begin{block}{Main Points}
    \begin{itemize}
      \item First main point.
      \item Second main point.
      \item Third main point.
    \end{itemize}
  \end{block}

  \begin{block}{Conclusion}
    In conclusion, the block environment is a versatile way to structure content in Beamer presentations.
  \end{block}
\end{frame}


\begin{frame}{Two-Column Slide Example}
  \begin{columns}[T] % Align columns at the top
    \begin{column}{0.47\textwidth}
      \textbf{Left Column}

      \begin{itemize}
        \item Item 1
        \item Item 2
        \item Item 3
      \end{itemize}
    \end{column}

    \begin{column}{0.47\textwidth}
      \textbf{Right Column}

      \begin{enumerate}
        \item First item
        \item Second item
        \item Third item
      \end{enumerate}
    \end{column}
  \end{columns}

  \vspace{3em} % Add some vertical space

  \textbf{Conclusion:}

  The `columns` environment is useful for creating multi-column layouts in Beamer presentations.
\end{frame}

\begin{frame}{Tabular Environment Example}
  \begin{tabular}{|c|c|c|}
    \hline
    \textbf{Column 1} & \textbf{Column 2} & \textbf{Column 3} \\
    \hline
    Row 1, Cell 1 & Row 1, Cell 2 & Row 1, Cell 3 \\
    \hline
    Row 2, Cell 1 & Row 2, Cell 2 & Row 2, Cell 3 \\
    \hline
    Row 3, Cell 1 & Row 3, Cell 2 & Row 3, Cell 3 \\
    \hline
  \end{tabular}

  \vspace{1em}

  \textbf{Conclusion:}

  The `tabular` environment is used to create tables in Beamer presentations. Adjust the formatting and content as needed.
\end{frame}


\begin{frame}{Pauses in presentation}
  \textbf{Introduction:}

  The \texttt{<>} structure after a command allows you to gradually reveal information in your presentation.

  \vspace{1em}

  \textbf{Example:}

  \begin{itemize}
    \item<1-> Start with an initial point.
    \item<2-> \pause Gradually reveal the next point.
    \item<3-> \pause Continue the sequence.
  \end{itemize}

  \vspace{1em}

  \textbf{Conclusion:}

  \begin{itemize}
    \item The \texttt{<>} structure command is useful for creating step-by-step presentations.
    \item It helps to keep the audience focused and engaged.
  \end{itemize}
\end{frame}



\begin{frame}{Using \textbackslash pause Command}
  \textbf{Introduction:}

  The \texttt{pause} command \pause can also be used.
  
  \vspace{3em}
  
  Other beamers commands such as \texttt{onslide, alt} can also be used.
\end{frame}

\begin{frame}
\frametitle{References}

You can use some of the macros I have supplied to provide references to other works. For example you can use \texttt{fcite} command\fcite{resnet}. You can use it multiple times\fcite{pascalvoc}. If you need to use the same citation again you should use \texttt{ffcite}\ffcite{resnet}.

\vspace{2em}

You should include the bib information in the \texttt{presentation.bib} file.
\end{frame}

\begin{frame}{Acknowledging Image Origin}
  To reference the source of an image you can use the \texttt{imfcite} command.

  \centering
  \includegraphics[width=0.5\textwidth]{example-image-a}\imfcite{resnet}

\end{frame}


\secframe{~}{Thank you for your attention!}

\appendix

\begin{frame}
\frametitle{Appendix}

You can add additional information in the appendix. These pages will not count towards the total number of pages.
\end{frame}

\secframe{~}{References}

\begin{frame}[allowframebreaks]
\frametitle{References}

\setbeamercolor{bibliography item}{fg=black}
\setbeamercolor*{bibliography entry title}{fg=black}
\setbeamercolor*{bibliography entry author}{fg=black}
\setbeamercolor*{bibliography entry location}{fg=black}
\setbeamercolor*{bibliography entry note}{fg=black}

\printbibliography
\end{frame}



\end{document}